\documentclass{ctexbook}
\usepackage{Note}
\usepackage{unicode-math}
\setCJKmainfont[
  BoldFont = 思源宋体 Bold,
  ItalicFont = 楷体,
]{宋体}
\begin{document}
\chapter{量子理论}
\section{量子力学的起源}
\subsection{能量的量子化}
\subsubsection{黑体辐射}
所有物体都会通过向外辐射电磁波的形式传递热量,这种传热方式称为\tbf{热辐射}.为了研究热辐射与物体温度的关系,假定一种理想的研究对象:
\begin{definition}[黑体]
    黑体是对一切入射的电磁辐射完全吸收而不发生反射或透射的物体.
\end{definition}
黑体的热辐射光谱仅与其温度有关,而与材质等无关.典型的黑体是一个开小孔的空腔.由经典理论可以推出黑体辐射的瑞利-金斯公式.
\begin{theorem}[瑞利-金斯公式]
    黑体辐射的辐射能量密度$u$与频率$\nu$,温度$T$的关系如下:
    \[u(\nu,T)=\dfrac{8\pi\nu^2kT}{c^3}\]
\end{theorem}
瑞利-金斯公式的推导过程如下:由于黑体辐射与黑体的体积,形状等都无关,因此考虑一个边长为$L$的立方体空腔.由于空腔的壁面上的电磁场必须为$0$,否则电磁波将穿透黑体,因此要求空腔内的电磁波形成驻波.因此其波长为
\[\lambda=\dfrac{2L}{n},\quad n=1,2,\cdots\]
我们用波数描述特定的波,有$k=\dfrac{2\pi}{\lambda}=\dfrac{n\pi}{L}$.在三维空间中,特定驻波的波矢可以按如下方式表述:
\[\vec{k}=\left(\dfrac{n_x\pi}{L},\dfrac{n_y\pi}{L},\dfrac{n_z\pi}{L}\right)\]
波矢的大小满足
\[\left|\vec{k}\right|^2=\dfrac{\pi^2}{L^2}n^2=\dfrac{\pi^2}{L^2}\left(n_x^2+n_y^2+n_z^2\right)\]
\subsection{波粒二象性}
\subsubsection{}
\section{波函数}
\subsection{波函数}
\subsection{薛定谔方程}
薛定谔方程是量子力学的基本假设之一;也可以启发式地从几个假设中推导出来:
\begin{enumerate}[label=\tbf{\alph*.}]
    \item 粒子的总能量$E$可以经典地表示为动能$T$和势能$V$的总和:
    \[E=T+V=\dfrac{p^2}{2m}+V\]
    \item 光子的能量与对应的电磁波的频率$f$成正比.假定这一事实对微观粒子成立,就有
    \[E=hf=\hbar\omega\]
    \item 德布罗意提出的假说表明微观粒子具有波粒二象性,其动量$p$和伴随的物质波波长$\lambda$有如下关系:
    \[p=\dfrac{h}{\lambda}=\hbar k\]
    其中$k=\dfrac{2\pi}{\lambda}$为波数.在空间中就有
    \[\vec{p}=\hbar\vec{k}\]
    \item 假定粒子的波函数是一个复平面波:
    \[\iPsi(x,t)=A\e^{\i(kx-\omega t)}\]
\end{enumerate}
\indent 于是根据\tbf{d.}可得波函数对时间$t$的导数和对位置$x$的二阶导数为
\[\dfrac{\p\Psi}{\p t}=A\e^{\i(kx-\omega t)}(-\i\omega)=-\i\omega\iPsi\]
\[\dfrac{\p^2\Psi}{\p x^2}=A\e^{\i(kx-\omega t)}(\i k)^2=-k^2\iPsi\]
根据\tbf{b.}可得
\[E\iPsi=\hbar\omega\iPsi=\i\hbar\dfrac{\p}{\p t}\iPsi\]
根据\tbf{c.}可得
\[p^2\iPsi=\hbar^2k^2\iPsi=-\hbar^2\dfrac{\p^2}{\p x^2}\iPsi\]
于是再根据\tbf{a.}可得
\[-\dfrac{\hbar^2}{2m}\dfrac{\p^2}{\p x^2}\iPsi+V\iPsi=\i\hbar\dfrac{\p}{\p t}\iPsi\]
推广到多维空间,则$\dfrac{\p^2}{\p x^2}$应当替换为Laplace算符.设哈密顿算符$\hat{H}=-\dfrac{\hbar^2}{2m}\nabla^2+V$,则有
\begin{theorem}[含时薛定谔方程]
    \[\hat{H}\iPsi=\i\hbar\dfrac{\p}{\p t}\iPsi,\quad\hat{H}=-\dfrac{\hbar^2}{2m}\nabla^2+V(x,t)\]
\end{theorem}
如果假定势函数$V$不含时,则$\hat{H}$也不含时.考虑一维情形的含时薛定谔方程(更高维度的也同理),做分离变量$\iPsi(x,t)=\psi(x)\phi(t)$,则有
\[-\dfrac{\hbar^2}{2m}\dfrac{\p^2}{\p x^2}\psi(x)\phi(t)+V(x)\psi(x)\phi(t)=\i\hbar\dfrac{\p}{\p t}\psi(x)\phi(t)\]
两边除以$\psi(x)\phi(t)$可得
\[-\dfrac{\hbar^2}{2m}\dfrac{1}{\psi(x)}\dfrac{\p^2}{\p x^2}\psi(x)+V(x)=\i\hbar\dfrac{1}{\phi(t)}\dfrac{\p}{\p t}\phi(t)\]
于是方程的两边分别是$x$和$t$的函数,因此它们都必须等于一个常数$E$(实际上就是能量).此时就有
\[\i\hbar\dfrac{\p}{\p t}\phi=E\phi,\quad-\dfrac{\hbar^2}{2m}\dfrac{\p^2}{\p x^2}\psi+V\psi=E\psi\]
第一个微分方程的解为
\[\phi(t)=\exp\left(-\dfrac{\i E}{\hbar}t\right)\]
而第二个方程即为
\begin{theorem}[定态薛定谔方程]
    \[\hat{H}\psi=E\psi,\quad\hat{H}=-\dfrac{\hbar^2}{2m}\nabla^2+V(x)\]
\end{theorem}
它的解依$V(x)$而定.
\section{算符}
\subsection{空间}
\begin{definition}[巴拿赫空间]
    完备的赋范向量空间是\tbf{巴拿赫空间}.
\end{definition}
\begin{definition}[希尔伯特空间]
    完备的内积空间是\tbf{希尔伯特空间}.
\end{definition}
\subsection{算符}
\begin{definition}[算符]
    算符是函数到函数的映射,即
    \[\hat{O}:f(x)\to g(x),\quad g(x)=\hat{O}f(x)\]
\end{definition}
\begin{definition}[伴随算符]
    定义在希尔伯特空间$\mathcal{H}$上的线性算符$\hat{A}$和$\hat{A}^\dagger$如果满足
    \[\avg{\hat{A}\vec{u},\vec{v}}=\avg{\vec{u},\hat{A}^\dagger\vec{v}},\forall\vec{u},\vec{v}\in\mathcal{H}\]
    则称$\hat{A}^\dagger$是$\hat{A}$的伴随算符.
\end{definition}
\begin{definition}[厄米算符]
    如果$\hat{A}^\dagger=\hat{A}$,则称$\hat{A}$是\tbf{自伴算符}(或\tbf{厄米算符}).
\end{definition}
\subsection{狄拉克符号与狄拉克$\symbf\delta$函数}
\subsubsection{狄拉克符号}
\begin{definition}[狄拉克符号]
    \tbf{狄拉克符号(Dirac notation)}是量子力学中广泛应用于描述量子态的一套标准符号系统.在这套系统中,每一个量子态都被描述为希尔伯特空间中的态矢量,定义为\tbf{右矢(ket)} $\ket{\psi}$;每一个右矢的共轭转置定义为其\tbf{左矢(bra)} $\bra{\psi}$.
\end{definition}
这样,态矢量$\ket{\psi}$与$\ket{\phi}$的内积可以表示为$\inprod{\psi}{\phi}$.
\subsubsection{狄拉克$\symbf\delta$函数}
有限维希尔伯特空间的正交归一基矢$\{\ket{e_i}\}_{i=1}^{N}$满足
\[\inprod{e_i}{e_j}=\delta_{ij}=\left\{\begin{array}{l}
    1,\quad i=j\\
    0,\quad i\neq j
\end{array}\right.\]
其中$\delta_{ij}$称作\tbf{克罗内克$\symbf\delta$符号}.对于无限维希尔伯特空间的正交归一基矢$\left\{\ket{x}\right\}$,其正交归一性质不能用离散的符号表示,因此定义如下函数:
\begin{definition}[狄拉克$\symbf\delta$函数]
    狄拉克$\delta$函数(通常简称为$\delta$函数)是满足以下条件的广义函数:
    \[\delta(x)\simeq\left\{\begin{array}{l}
        +\infty,\quad x=0\\
        0,\quad x\neq 0
    \end{array}\right.,\quad\int_{-\infty}^{+\infty}\delta(x)\d x=1\]
\end{definition}
$\delta$函数的常见表达式有下面几种:
\[\delta(x)=\dfrac{1}{2\pi}\int_{-\infty}^{+\infty}\e^{\i kx}\d k,\quad\delta(x)=\lim_{\ep\to0^+}\dfrac{1}{2\sqrt{\pi\ep}}\exp\left(-\dfrac{x^2}{4\ep}\right),\quad\delta(x)=\lim_{\ep\to0}\dfrac{\sin(x/\ep)}{\pi x}\]
从第一个形式可以看出$\delta$函数和常数函数互为傅里叶变换的关系;从第二个形式可以看出$\delta$函数可以视作方差无穷小的高斯分布.$\delta$函数有如下重要性质:
\begin{theorem}
    对于任意在$x=0$处连续的函数$f(x)$,都有
    \[\int_{-\infty}^{+\infty}f(x)\delta(x)\d x=f(0)\]
\end{theorem}
\indent 引入狄拉克函数后,无穷维线性空间中的基矢的正交归一性即可由
\[\inprod{x}{x'}=\delta(x-x')\]
描述.
\subsubsection{基组展开}
对于希尔伯特空间中的离散基向量组$\left\{u_j\right\}_{j=1}^{\infty}$,考虑某一向量$\ket{v}$表示为这组基的线性组合:
\[\ket{v}=\sum_{j=1}^{\infty}c_j\ket{u_j}\]
而
\[\inprod{u_j}{v}=\bra{u_j}\left(\sum_{k=1}^{\infty}c_k\ket{u_k}\right)=\sum_{k=1}^{\infty}c_k\inprod{u_j}{u_k}=\sum_{k=1}^{\infty}c_k\delta_{kj}=c_j\]
于是有
\begin{theorem}[离散基组下的展开]
    \[\ket{v}=\sum_{j=1}^{\infty}\inprod{u_j}{v}\ket{u_j}\]
\end{theorem}
考虑到
\[\ket{v}=\sum_{j=1}^{\infty}\inprod{u_j}{v}=\sum_{j=1}^{\infty}\ket{u_j}\inprod{u_j}{v}=\sum_{j=1}^{\infty}\left(\ket{u_j}\bra{u_j}\right)\ket{v}=\left(\sum_{j=1}^{\infty}\ket{u_j}\bra{u_j}\right)\ket{v}\]
于是可得
\begin{theorem}[离散基组下的恒等算符]
    \[\hat{I}=\sum_{j=1}^{\infty}\ket{u_j}\bra{u_j}\]
\end{theorem}
在很多推导中都可以用到这一等式.\\
\indent 同样地,在连续基$\ket{x}$下态$\ket{v}$的展开为
\[\ket{v}=\int_{-\infty}^{+\infty}f(x)\ket{x}\d x\]
同样地有
\[\inprod{x}{v}=\bra{x}\left(\int_{-\infty}^{+\infty}f(x')\ket{x'}\d x'\right)=\int_{-\infty}^{+\infty}f(x')\inprod{x}{x'}\d x'=\int_{-\infty}^{+\infty}f(x')\delta(x-x')\d x'=f(x)\]
于是类似地有
\begin{theorem}[连续基组下的展开和恒等算符]
    \[\ket{v}=\int_{-\infty}^{+\infty}\inprod{x}{v}\ket{x}\d x,\quad\hat{I}=\int_{-\infty}^{+\infty}\ket{x}\bra{x}\d x\]
\end{theorem}
\subsubsection{表象}
\indent 对于希尔伯特空间中的量子态$\ket{\psi}$,我们可以用不同的基对其进行描述:
\begin{enumerate}[label=\tbf{\alph*.}]
    \item \tbf{位置空间(实空间)}.位置空间的基矢为$\ket{x}$,表示粒子位于$x$位置的态. $\ket{\psi}$在此连续基下展开的形式为
    \[\ket{\psi}=\int\inprod{x}{\psi}\ket{x}\d x\]
    记$\psi(x)=\inprod{x}{\psi}$即该态在位置空间中的表示.大部分时候,我们研究的(空间)波函数都是表示在位置空间中的,它们总是具有$\psi(x)$的形式.
    \item \tbf{动量空间}.动量空间的基矢为$\ket{k}$,表示粒子动量为$\ket{k}$的态. $\ket{\psi}$在此连续基下展开的形式为
    \[\ket{\psi}=\int\inprod{k}{\psi}\ket{k}\d k\]
    记$\phi(k)=\inprod{k}{\psi}$即该态在位置空间中的表示.
    \item \tbf{函数空间}.考虑某一算符对应的本征态构成的基$\{\varphi_j\}_{j=1}^{\infty}$,它们都对应于某种运动模式. $\ket{\psi}$在此基组下展开的形式为
    \[\ket{\psi}=\sum_{j=1}^{\infty}\inprod{\varphi_j}{\psi}\ket{\varphi_j}\]
\end{enumerate}
所谓位置空间,动量空间和函数空间,都是我们讨论的希尔伯特空间在某一组基下的表现形式,即所谓\tbf{表象}.\\
\indent 比较重要的是同一态在位置空间中的表示$\psi(x)$和在动量空间中的表示$\phi(k)$之间的关系.在推导之前,我们需要\textit{假定}\footnote{可以从后面的不对易关系推导出来.}动量算符$\hat{p}$在位置空间中的表现形式为
\[\hat{p}=-\i\hbar\dfrac{\p}{\p x}\]
那么动量算符的本征态$\ket{k}$在位置空间中的表示$\psi_{k}(x)$满足
\[k\psi_{k}(x)=\hat{p}\psi_k(x)=-\i\hbar\dfrac{\p}{\p x}\psi_k(x)\]
这一微分方程的解为$\psi_{k}(x)=C\e^{\i kx}$,归一化即可得
\[\inprod{x}{k}=\psi_k(x)=\dfrac{1}{\sqrt{2\pi}}\e^{\i kx}\]
于是
\[\phi(k)=\inprod{k}{\psi}=\bra{k}\hat{I}\ket{\psi}=\bra{k}\left(\int_{-\infty}^{+\infty}\ket{x}\bra{x}\d x\right)\ket{\psi}=\int_{-\infty}^{+\infty}\inprod{k}{x}\inprod{x}{\psi}\d x=\dfrac{1}{\sqrt{2\pi}}\int_{-\infty}^{+\infty}\psi(x)\e^{-\i kx}\d x\]
这就证明了两者互成傅里叶变换关系.
\subsection{对易关系和态叠加原理}
\subsubsection{对易关系}
为了衡量两个算符的不可交换性,定义
\begin{theorem}[对易子和对易算符]
    考虑两个希尔伯特空间中的算符$\hat{A}$和$\hat{B}$,定义它们的\tbf{对易子}
    \[[\hat{A},\hat{B}]=\hat{A}\hat{B}-\hat{B}\hat{A}\]
    如果$[\hat{A},\hat{B}]=0$,则称$\hat{A}$和$\hat{B}$\tbf{对易}.
\end{theorem}
对易算符的本征态有如下关系:
\begin{theorem}
    希尔伯特空间中的算符$\hat{A}$和$\hat{B}$对易当且仅当它们具有完全相同的本征态(更严格的,具有相同的本征子空间).
\end{theorem}
\subsubsection{态叠加原理}
前面已经说过,态矢量$\ket{\psi}$可以用基组展开的形式表示:
\[\ket\psi=\sum_{j=1}^{\infty}c_j\ket{\varphi_j}\]
量子力学讲上述线性组合解释为\tbf{态的叠加}.而对于本征态$\varphi_j$的物理意义,量子力学给出了如下的基本假设:
\begin{theorem}[可观测量与厄米算符]
    量子力学中的可观测量都对应于一个厄米算符.
\end{theorem}
\begin{theorem}[态叠加原理]
    考虑测量态$\ket{\psi}$的某一可观测量$H$,其对应的厄米算符为$\ket{H}$,则
    \begin{enumerate}[label=\tbf{\alph*.}]
        \item 如果$\ket{\psi}$是$\hat{H}$的本征态,那么测量结果是$\ket{\psi}$对应的本征值.
        \item 如果$\ket{\psi}$是$\hat{H}$的本征态组成的叠加态,即
        \[\ket\psi=\sum_{j=1}^{\infty}c_j\ket{\varphi_j}\]
        那么测量结果是某个本征态$\ket{\varphi_j}$对应的本征值$\lambda_j$,出现该结果的概率是$\left|c_j\right|^2$.得到这一结果的过程称作\tbf{坍缩}.
    \end{enumerate}
\end{theorem}
根据态叠加原理可得
\begin{theorem}[对易关系与测量]
    对易算符描述的可观测量可以同时准确描述;不对易算符描述的可观测量不可以同时描述.
\end{theorem}
\begin{theorem}[不确定原理]
    设可观测量$A$和$B$对应的算符为$\hat{A}$和$\hat{B}$.定义算符的标准差
    \[\sigma_{O}=\sqrt{\avg{\hat{O}^2}-\hat{O}^2}\]
    则有
    \[\sigma_A^2\sigma_B^2\geq\left(\dfrac{1}{2\i}\avg{[\hat{A},\hat{B}]}\right)^2\]
\end{theorem}
\section{平动,振动与转动}
\subsection{平动}
\subsubsection{一维自由粒子}
一维自由粒子的势函数$V(x)=0$.定态薛定谔方程为
\[-\dfrac{\hbar^2}{2m}\dfrac{\d^2\psi(x)}{\d x^2}=E\psi(x)\]
该方程的通解为
\[\psi_k(x)=A\e^{\i kx}+B\e^{-\i kx},\quad E_k=\dfrac{k^2\hbar^2}{2m}\]
可以看出,一维自由粒子的能量是\tbf{连续取值}的,这是因为没有势函数的限制.它的波函数$\iPsi_k(x)$也无法归一化,是两列平面波的叠加.
\subsubsection{一维无限深势阱}
考虑一个无限深的一维势阱,其势函数为
\[V(x)=\left\{\begin{array}{l}
    0,\quad 0\leq x\leq L\\
    +\infty,\quad\text{otherwise}
\end{array}\right.\]
在势阱以外, $\psi(x)=0$(也即在那里发现粒子的几率为$0$),而在势阱内定态薛定谔方程为
\[-\dfrac{\hbar^2}{2m}\dfrac{\di^2\psi}{\di x^2}=E\psi\]
我们已经知道$E\geq V_{\min}=0$,于是则有
\[\dfrac{\di^2\psi}{\di x^2}=-k^2\psi,\quad k=\dfrac{\sqrt{2mE}}{\hbar}\]
上述微分方程的通解为
\[\psi(x)=A\sin kx+B\cos kx\]
$\psi(x)$的连续性要求
\[\psi(0)=\psi(L)=0\]
于是
\[\psi(0)=B=0,\quad \psi(L)=A\sin kx=0\]
如果$A=0$,那么$\psi(x)=0$,没有物理意义.这意味着$\sin ka=0$,于是
\[k_n=\dfrac{n\pi}{L},\quad n=1,2,\cdots\]
然后根据归一化条件可得
\[\int_0^a\left|A\right|^2\sin^2(kx)\di x=\left|A\right|^2\dfrac{L}{2}=1\]
取$A=\sqrt{\dfrac{2}{L}}$,则无限深方势阱内的定态薛定谔方程的解为
\[\psi_n(x)=\sqrt{\dfrac{2}{L}}\sin\left(\dfrac{n\pi}{L}x\right),\quad E_n=\dfrac{n^2\pi^2\hbar^2}{2mL^2}=\dfrac{n^2h^2}{8mL^2},\quad n=1,2,\cdots\]
与一维自由粒子对比,可以发现势阱的限制使得能量量子化.\\
\indent 更高维度的无限深方势阱可以通过分离变量的方法转化为各个维度上的无限深方势阱问题,并采用类似的办法求解.\\
\indent 无限深势阱可以应用于描述密闭容器内稀薄气体分子的运动,定性地预测长链共轭体系的激发能量等.
\subsubsection{一维有限宽有限深势垒}
考虑以下势函数
\[V(x)=\left\{\begin{array}{l}
    V_0,\quad0\leq x\leq W\\
    0,\quad\text{otherwise}
\end{array}\right.\]
这一体系的本征态分多种情况,这里考虑势垒左边是平面波的情形.首先假定$E<V_0$,即能量低于势垒的情形.分三段考虑波函数:
\[\psi(x)=\left\{\begin{array}{l}
    A\e^{\i kx}+B\e^{-\i kx},\quad x\leq0\\
    C\e^{\kappa x}+D\e^{-\kappa x},\quad 0<x<W\\
    A'\e^{\i kx},\quad x\geq W
\end{array}\right.\quad\text{where}\quad k=\dfrac{\sqrt{2mE}}{\hbar},\quad \kappa=\dfrac{\sqrt{2m(V_0-E)}}{\hbar}\]
根据$\psi(x)$和其一阶导数$\dfrac{\d\psi(x)}{\d x}$的连续性可得
\[x=0:\quad A+B=C+D,\quad \i kA-\i kB=\kappa C-\kappa D\]
\[x=W:\quad C\e^{\kappa W}+D\e^{-\kappa W}=A'\e^{\i k W},\quad \kappa C\e^{\kappa W}+\kappa D\e^{-\kappa W}=\i kA'\e^{\i k W}\]
一共有四个方程和五个待定系数,可以导出它们的比例关系.定义\tbf{透射概率} $T=\dfrac{|A'|^2}{|A|^2}$,整理可得
\[T=\left[1+\dfrac{(\e^{\kappa W}-\e^{-\kappa W})^2}{16\ep(1-\ep)}\right]^{-1},\quad \ep=\dfrac{E}{V_0}\]
对于宽势垒有$\kappa W\gg1$,此时
\[T=16\ep(1-\ep)\e^{-2\kappa W}\]
\subsection{振动}
简谐振子的势函数为
\[V(x)=\dfrac12 k_\text{f}x^2\]
求解可得
\[E_v=\left(v+\dfrac12\right)\hbar\omega,\quad\omega=\sqrt{\dfrac{k_\text{f}}{m}}\]
其中$v$是\tbf{振动量子数}. \tbf{振动零点能}$E_0=\dfrac12\hbar\omega$.振动波函数为
\[\psi_v(x)=N_vH_v(y)\e6^{-y^2/2},\quad y=\dfrac{x}{\alpha},\quad\alpha=\left(\dfrac{\hbar^2}{mk_\text{f}}\right)^{1/4}\]
其中$H_v$是厄米多项式,满足
\[H''_v-2yH'_v+2vH_v=0\]
振子的基态波函数是高斯函数.
\subsection{转动}
\subsubsection{二维转动}
考虑限制在二维平面上半径为$r$的圆周上运动的粒子,其定态薛定谔方程为
\[-\dfrac{\hbar^2}{2m}\nabla^2\psi=E\psi\]
而
\[\nabla^2=\dfrac{\p^2}{\p x^2}+\dfrac{\p^2}{\p y^2}=\dfrac{\p}{\p r^2}+\dfrac1r\dfrac{\p}{\p r}+\dfrac{1}{r^2}\dfrac{\p^2}{\p\phi^2}\]
由于$r$固定为常数,因此$\psi$只是$\phi$的函数,从而
\[-\dfrac{\hbar^2}{2m r^2}\dfrac{\d^2\psi}{\d\phi^2}=E\psi\]
根据周期性可得$\psi(\phi)=\psi(\phi+2\pi)$,于是解得
\[\psi_{m_l}(\phi)=\dfrac{\e^{\i m_l\phi}}{\sqrt{2\pi}},\quad E_{m_l}=\dfrac{m_l^2\hbar^2}{2I},\quad m_l=0,\pm1,\pm2,\cdots\]
其中转动惯量$I=mr^2$.正是周期性的边界条件导致体系的量子化.\\
\indent 现在考虑角动量$\vec{I}$.在经典力学中有$\vec{l}=\vec{r}\times\vec{p}$,由于粒子在$xy$平面上运动,则有$l_z=xp_y-yp_x$.将$p_x$和$p_y$用量子力学的动量算符代入可得角动量算符为
\[\hat{l}_z=\dfrac{\hbar}{\i}\left(x\dfrac{\p}{\p y}-y\dfrac{\p}{\p x}\right)=\dfrac{\hbar}{\i}\dfrac{\d}{\d\phi}\]
于是对于体系的本征函数$\psi_{m_l}(\phi)$,其角动量本征值为
\[\dfrac{\hat{l}_z\psi_{m_l}(\phi)}{\psi_{m_l}(\phi)}=\dfrac{\dfrac{\hbar}{\i}\dfrac{\d}{\d\phi}\e^{\i m_l\phi}}{\e^{\i m_l\phi}}=m_l\hbar\]
\indent 可以看出,位置$\phi$和角动量的本征态不同,因此两者不能同时测准.而动能与角动量共用一组本征波函数.
\subsubsection{三维转动}
同样地,考虑限制在三维空间上半径为$r$的球面上运动的粒子,其定态薛定谔方程为
\[-\dfrac{\hbar^2}{2m}\nabla^2\psi=E\psi,\quad\nabla^2=\dfrac{1}{r^2}\dfrac{\p}{\p r}\left(r^2\dfrac{\p}{\p r}\right)+\dfrac{1}{r^2\sin\theta}\dfrac{\p}{\p\theta}\left(\sin\theta\dfrac{\p}{\p\theta}\right)+\dfrac{1}{r^2\sin^2\theta}\dfrac{\p^2}{\p\phi^2}\]
同样地,考虑将$\nabla^2$分为与$r$有关和无关的两部分
\[\nabla^2=\dfrac{1}{r^2}\dfrac{\p}{\p r}\left(r^2\dfrac{\p}{\p r}\right)+\dfrac{1}{r^2}\underbrace{\left[\dfrac{1}{\sin\theta}\dfrac{\p}{\p\theta}\left(\sin\theta\dfrac{\p}{\p\theta}\right)+\dfrac{1}{\sin^2\theta}\dfrac{\p^2}{\p\phi^2}\right]}_{\Lambda^2}\]
其中$\Lambda^2$是勒让德算子.于是$\psi$是$\theta,\phi$的函数,薛定谔方程改写为
\[-\dfrac{\hbar^2}{2I}\Lambda^2\psi=E\psi\]
这一微分方程的解是\tbf{球谐函数}$Y_{l,m_l}$,需要用两个量子数$l$(称作\tbf{角量子数}), $m_l$(称作\tbf{磁量子数})描述体系:
\[E_{l,m_l}=l(l+1)\dfrac{\hbar^2}{2I},\quad l=0,1,2,\cdots,\quad m_l=0,\pm1,\cdots,\pm l\]
\indent 现在考虑角动量与角动量$z$分量的关系,有
\begin{theorem}[角动量与角动量分量的对易关系]
    \[\left[\hat{l^2},\hat{l}_z\right]=0\]
\end{theorem}
因此两者具有相同的本征态.实际上,球谐函数总是可以写成
\[Y_{l,m_l}=f(\theta)\e^{\i m_l\phi}\]
的形式,因此它也是$\hat{l}_z$的本征态.而不难证明
\begin{theorem}[角动量分量之间的对易关系]
    \[\left[\hat{l}_x,\hat{l}_y\right]=\i\hbar\hat{l}_z,\quad\left[\hat{l}_y,\hat{l}_z\right]=\i\hbar\hat{l}_x,\quad\left[\hat{l}_z,\hat{l}_x\right]=\i\hbar\hat{l}_y\]    
\end{theorem}
因此角动量的各分量之间不对易,它们不能同时测准.
\end{document}